% !PRINTED 598
\weberpage{598}
%<<CONT
eine \gesperrt{metacyklische Gleichung} nennen. Die cyklischen
Gleichungen selbst sind als specieller Fall darunter mit enthalten,
und nach dem im \S.\,175 Bemerkten sind die metacyklischen
Gleichungen dieselben, wie die durch Radicale lösbaren Gleichungen. Ist $P$ die Gruppe einer solchen Gleichung, so muss sie
nach dem vorigen Paragraphen einen Normaltheiler von Primzahlindex $j_1$, den wir jetzt mit $P_1$ bezeichnen wollen, besitzen.
Besteht $P_1$ aus der einzigen identischen Permutation, so ist $P$
selbst cyklisch und von Primzahlgrad. Ist $P_1$ nicht die Einheitsgruppe, so muss $P_1$ wieder einen Normaltheiler $P_2$ von Primzahlindex $j_2$ enthalten u. s. f., bis wir endlich zur Einheitsgruppe
gelangen.

Dass es auch die für eine metacyklische Gleichung ausreichende Bedingung ist, wenn ihre Gruppe $P$ diese Zusammensetzung hat, ergiebt sich aus dem vorigen Paragraphen. Wir
sprechen also den Satz aus:

\begin{quote}
\noindent III.\ \gesperrt{Die nothwendige und hinreichende Bedingung für
eine metacyklische Gleichung ist die, dass es
eine Reihe von Gruppen}

\[
P, P_1, P_2, P_3 \ldots
\]

\gesperrt{giebt, deren erste die Galois'sche Gruppe der
Gleichung, deren letzte die Einheitsgruppe ist,
von denen jede folgende ein normaler Theiler
der nächst vorangehenden von Primzahlindex ist.}
\end{quote}

Hiernach nennen wir eine Permutationsgruppe $P$, die diese
Eigenschaft hat, zu der sich also eine Kette von Gruppen

\[
P, P_1, P_2 \ldots P_{\mu-1}, 1
\]

so bestimmen lässt, dass jedes Glied Normaltheiler des vorangehenden von Primzahlindex ist, eine \gesperrt{metacyklische
Gruppe}\footnote{Der Ausdruck „metacyklische Gruppen“ ist zuerst von \gesperrt{Kronecker},
wenn auch in beschränkterem Sinne gebraucht. Ich möchte hier diese
leichte Verallgemeinerung eines schon bekannten Ausdruckes an Stelle des
von \gesperrt{Frobenius} und \gesperrt{Hölder} benutzten Ausdruckes „auflösbare Gruppen“
vorschlagen.}.

Wir haben hier die Bedingung für die \gesperrt{vollständige} Auflösbarkeit einer Gleichung durch eine Kette von cyklischen
Gleichungen erhalten. Es handelt sich aber noch um die Frage,
%CONT>>
